% 第2章 熱平衡宇宙の熱史（仕様 01_textbook_spec.md §第2章）
\section{熱平衡宇宙の熱史}
\label{ch:ch02_thermal}

\paragraph{この章で導くこと（入力→出力）}
平衡統計力学を膨張宇宙に適用し、$T\propto1/a$ の黒体放射、ニュートリノ脱結合と
$T_\nu/T_\gamma=(4/11)^{1/3}$（契約 S2.1）、相互作用率 vs 膨張率による脱結合基準を
確立する。CMB が黒体である理由をここで完結させる。

\subsection{分布関数と熱力学量}
平衡分布 $f(\mathbf p)=[\exp((E-\mu)/k_BT)\mp1]^{-1}$（上：Bose、下：Fermi）から、
内部自由度 $g$ の粒子の数密度・エネルギー密度を
\begin{equation}
  n=\frac{g}{(2\pi\hbar)^3}\!\int\! f\,d^3p,\qquad
  \rho=\frac{g}{(2\pi\hbar)^3}\!\int\! E(p)\,f\,d^3p,
\end{equation}
で得る。相対論的極限 $E=pc$、$\mu=0$ では $\int_0^\infty\frac{p^3dp}{e^{pc/k_BT}\mp1}$ を
無次元化（$u=pc/k_BT$）して $\int u^3/(e^u\mp1)du$ を使うと
\begin{equation}
  \rho_{\rm boson}=\frac{\pi^2}{30}\frac{(k_BT)^4}{(\hbar c)^3}\,g,\qquad
  \rho_{\rm fermion}=\frac{7}{8}\,\rho_{\rm boson}\big|_{g\to g},
  \label{eq:rho_rel}
\end{equation}
を得る（Fermi の $7/8$ は $\int u^3/(e^u+1)=\tfrac78\int u^3/(e^u-1)$ から）。これらを
$\rho=\tfrac{\pi^2}{30}g_*(k_BT)^4/(\hbar c)^3$ とまとめ $g_*$ を有効自由度と呼ぶ。
非相対論極限では $\rho\propto(k_BT)^{3/2}e^{-mc^2/k_BT}$ と Boltzmann 抑制される。

\subsection{膨張宇宙での平衡維持条件}
相互作用率 $\Gamma=n\langle\sigma v\rangle$ と膨張率 $H$ の比較で脱結合を判定する。
$\Gamma\gtrsim H$ で平衡、$\Gamma\lesssim H$ で脱結合（フリーズアウト）。

\begin{intuition}
「反応が1膨張時間に1回以上起きるか」が平衡維持の目安である。膨張が反応に
追いつくと粒子は相互作用を失い、共動数密度が凍結する。
\end{intuition}

\subsection{エントロピー保存と $T(a)$}
共動エントロピー $s\,a^3=$ const から $g_{*s}T^3a^3=$ const。$g_{*s}$ が一定の間は
$T\propto1/a$。

\subsection{ニュートリノ脱結合（導出契約 S2.1）}
ニュートリノは $T\sim1$~MeV で脱結合し、以後 $T_\nu\propto a^{-1}$ で独立に冷える。
その直後 $T\sim m_ec^2$ で $e^+e^-$ が消滅し、そのエントロピーは\emph{光子のみ}を加熱
する。エントロピー自由度 $g_{*s}$ は、消滅前は光子($2$)＋$e^\pm$($\tfrac78\times4$)で
$g_{*s}^{\rm before}=2+\tfrac72=\tfrac{11}{2}$、消滅後は光子のみで $g_{*s}^{\rm after}=2$。
光子側のエントロピー保存 $g_{*s}T_\gamma^3a^3=$ const より
\begin{equation}
  \Big(\frac{T_\gamma}{T_\nu}\Big)^3=\frac{g_{*s}^{\rm before}}{g_{*s}^{\rm after}}
  =\frac{11/2}{2}=\frac{11}{4}
  \;\Longrightarrow\;
  \frac{T_\nu}{T_\gamma}=\left(\frac{4}{11}\right)^{1/3}\approx0.714.
  \label{eq:Tnu}
\end{equation}
（$T_\nu$ は $e^\pm$ 加熱を受けないので比が生じる。）相対論的1種あたりのエネルギー密度は
式\eqref{eq:rho_rel} の $7/8$ と $(T_\nu/T_\gamma)^4=(4/11)^{4/3}$ から
\begin{equation}
  \Omega_{\nu0}=N_{\rm eff}\,\frac{7}{8}\left(\frac{4}{11}\right)^{4/3}\Omega_{\gamma0}
\end{equation}
（質量ゼロ）。これが背景\eqref{eq:Hp} の放射項と第\ref{ch:ch07_initial}章の $f_\nu$ に入る。

\begin{advanced}
$N_{\rm eff}=3.046$ の $0.046$ は、脱結合が瞬間的でなく $e^+e^-$ 消滅と一部
重なるための補正である。本書の主筋では質量ゼロ・$N_{\rm eff}=3.046$ を用いる。
\end{advanced}

\subsection{黒体性の維持（導出契約 S2.3）}
Liouville 方程式により、無衝突自由膨張下で位相空間分布 $f(p)$ は不変に保たれる。
プランク分布は $f=[\exp(pc/k_BT)-1]^{-1}$ の形を $T\propto1/a$ のもとで保つため、
CMB は冷えても黒体であり続ける（FIRAS の観測と整合）。

\begin{numreality}
光子のエネルギー密度から $\Omega_{\gamma0}=\dfrac{a_{\rm rad}T_{\rm CMB}^4/c^2}{\rho_{\rm crit,0}}$、
$a_{\rm rad}=\pi^2k_B^4/(15\hbar^3c^3)$。$T_{\rm CMB}=2.7255$~K を入れると
$\Omega_{\gamma0}h^2\approx2.47\times10^{-5}$（契約 S2.2、\texttt{cmbcore.params} と一致）。
\end{numreality}

\subsection{トムソン散乱とコンプトン散乱}
微分断面積 $\frac{d\sigma_T}{d\Omega}=\frac{3\sigma_T}{16\pi}(1+\cos^2\theta)$。
再結合期では $h\nu\sim k_BT_*\sim0.3$~eV $\ll m_ec^2=511$~keV ゆえエネルギー移行は
$O(h\nu/m_ec^2)\sim10^{-6}$ に抑制され、散乱は光子の\emph{方向のみ}を変える。これが
第\ref{ch:ch06_boltzmann}章のボルツマン衝突項（Thomson 項）の物理的基礎となる。

\paragraph{まとめ}
熱史の登場人物（計量・光子・ニュートリノ・バリオン・CDM・$\Lambda$）が確定し、
光子は黒体、再結合期の散乱は方向変化のみ、という第II部への前提が揃った。

\paragraph{演習}
(1) $\Omega_{\gamma0}h^2$ を $T_{\rm CMB}$ から数値導出せよ（契約 S2.2）。
(2) $T_\nu/T_\gamma$ をエントロピー保存から導け。
(3) トムソン極限の妥当性を再結合期の典型光子エネルギーで数値確認せよ。
